\section{Discussion}
\subsection{Principal findings and relation to prior work}
The Conformer pipeline improved held-out accuracy from 0.14--0.27 (CNN-only, CNN-LSTM, and attention-only baselines) to 0.60 while retaining a compact architecture and reproducible preprocessing chain. These gains highlight the benefit of combining convolutional feature extraction with attention-based context when working with Morlet scalograms, complementing recent reports on residual-attention classifiers and interpretable architectures for biosignals.\cite{rahman2024residual,rowlandson2024bridging} At the same time, the macro-F1 ceiling of 0.26 underscores how severe class imbalance and minority morphologies constrain arrhythmia screening pipelines, echoing wearable-focused studies that advocate context-aware feature learning to stabilise predictions under noisy acquisition.\cite{thapa2025generalizable,tison2018passive} Publishing this conservative, fully scripted baseline reduces the risk of optimistic reporting and establishes a clear starting point for augmentation-focused follow-on work.

To situate these numbers within the broader MIT-BIH literature, \autoref{tab:literature-context} compiles representative multi-class arrhythmia classifiers. Many headline accuracies above 90\% are obtained under beat-wise or patient-specific splits---for example, the deep CNN of Acharya et al.\cite{acharya2017deep}---where individual beats from the same subject may appear in both training and test sets. In contrast, the present work adheres to an inter-patient, record-level split with the AAMI five-class mapping, which yields lower macro-F1 but more conservative estimates of generalisation to unseen patients.

\subsection{Implications for telecardiology workflows}
Operational deployment will hinge on pairing the released ROC, precision--recall, and confusion-matrix diagnostics with clinician-defined thresholds so that minority-class misses trigger manual review. Because supraventricular and fusion rhythms remain poorly detected, the present model should be used only as a triage aid or benchmarking tool and not as a stand-alone arrhythmia detector. Although this work does not yet include temperature scaling or workload simulations, the open \texttt{Research\_Runs} artifacts expose raw probabilities, predictions, and figure notebooks for downstream calibration, aligning with deployments where reproducibility and auditability drive adoption.\cite{rowlandson2024bridging,spagna2024telecardiology,tison2018passive} The modular preprocessing pipeline also facilitates adaptation to longer-duration wearables or implantable recorders without rewriting the training stack.

\subsection{Limitations and future directions}
Residual errors on fusion and supraventricular beats reveal the need for targeted augmentation, alternative loss formulations, or clinician-in-the-loop relabeling to mitigate scarcity even in curated repositories. Expanding the corpus with synthetic signals or cross-dataset fine-tuning may help but risks distribution shift, as highlighted by investigations into transferability across imaging and motion-tracking modalities.\cite{thapa2025generalizable,jenifer2024yolo} In addition, the bootstrap confidence intervals reported for accuracy, macro-F1, and AUCs are estimated at the beat level and therefore treat individual beats as independent; they quantify internal variability on the held-out test cohort rather than population-level uncertainty across patients. Future work will pair Conformer feature extraction with probabilistic calibration layers, including simple post-hoc schemes such as class-wise temperature scaling, and federated training schemes to satisfy privacy constraints in multi-center telehealth deployments. We also did not yet perform robustness experiments in the signal or scalogram domain (e.g., additive noise on the ECG or Morlet images); these analyses remain natural extensions to the present probability-space perturbations and could help clarify how acquisition artefacts propagate into clinical decision thresholds. Collaborations with pediatric cardiology programs could examine age-specific morphology effects similar to those reported for pubertal ECG analytics.\cite{spagna2024telecardiology}

\subsection{Ethical considerations}
The study exclusively utilized de-identified public datasets, and all analytic code, trained weights, and derived outputs are released alongside this manuscript to support independent verification. No human subjects were recruited, obviating the need for institutional review, and the analytical workflow complies with the project’s documented ethics and data availability statements.

\begin{table*}[t]
  \centering
  \caption{Representative multi-class arrhythmia classifiers on the MIT-BIH Arrhythmia Database. Reported metrics are quoted from the original papers where available; split definitions highlight that many high headline accuracies rely on beat-wise or patient-specific evaluation rather than fully inter-patient, five-class AAMI mapping.}
  \label{tab:literature-context}
  \resizebox{\textwidth}{!}{%
  \begin{tabular}{lllll}
    \toprule
    Paper & Dataset / split & Classes & Acc./Macro-F1 & Notes \\
    \midrule
    This work (Morlet--Conformer) & MIT-BIH, 45 records, 38 train/val / 7 test, inter-patient & 5 (AAMI N/S/V/F/Q) & 0.60 / 0.26 & Record-level split; strictly held-out test subjects; beat-level bootstrap CIs. \\
    Acharya et al.\ (2017)\cite{acharya2017deep} & MIT-BIH, beat-wise random train/test split & 5 heartbeat categories & $\approx 0.94$ / -- & High beat-wise accuracy; macro-F1 not reported; not strictly inter-patient. \\
    Kiranyaz et al.\ (2015)\cite{kiranyaz2015personalized} & MIT-BIH, patient-specific adaptive training & Multiple arrhythmia types & -- & Real-time 1-D CNN with subject-specific adaptation; performance not directly comparable to fixed inter-patient split. \\
    \bottomrule
  \end{tabular}
  }
\end{table*}

